%----------------------------------------------------------------------------
\chapter{\bevezetes}
%----------------------------------------------------------------------------

Ahogy minden projekt, ez is egy problémával indult. Minden csoport vagy közösség eseményein előfordul a szükség arra, hogy a közösen megvásárolt termékeket a közösségen belül elfogyaszthassák a tagok, és becsületkassza alapon bedobják az összegét az elvett terméknek. Ezt legtöbbször egy papíron vezetik strigulázva a nevek mellett vagy jobb esetben egy táblázatkezelő szoftverbe vezetik, hogy ki mit fogyasztott a közösből. Bár ezek megoldják az alapvető problémát, de nem tökéletes. Nem lehet velük visszamenőleg követni, hogy ki, mikor, mit csinált, mennyi az egyenlege, stb\ldots

Ezért aztán jött az ötlet, hogy mérnökinformatikusként nehogy már papíron vezessük ezeket a kiadásokat, hanem legyen erre egy szoftver, ami ezt követni, és mindenféle okos funkcióval megkönnyíti és megoldja az említett problémát. Így született meg egy évvel ezelőtt a KasszaSCH ötlete.

Természetesen, nem ugrottam fejest rögtön a kódolásba, hiszen egy szoftver megtervezése előtt elő kell állni a követelményekkel, hogy mire lenne szükségük a felhasználóknak, kutatást végezni, hogy van-e erre már egy jól bevált program, vagy annyira speciálisak az igények, hogy egyszerűbb egy sajátot készíteni. Sok emberrel való konzultáció után megszülettek az igények, hogy mit is kéne tudjon a jövőbeli alkalmazás. Ezen belül megállapítottuk, mik lennének a legfontosabb funkciók, és mik a "jó lenne, ha ezt is tudná" kérések.

Fél évvel később, a stakeholderekkel való hosszas beszélgetés, a követelmények elemzése és megfogalmazása, valamint időhiány után elkezdődött a KasszaSCH fejlesztése.

Dolgozatomban ennek a rendszernek a tervezését és megvalósítását fogom bemutatni.
